\section{Corporate restructuring}

\begin{qfmethod}[Evaluating corporate restructurings: types and motives, the initial evaluation, investment actions (acquisitions, joint ventures, stakes), divestment actions (sales, spin-offs, split-offs, carve-outs) and cost and balance-sheet restructurings]\label{meth:restructuring-evaluation}
\qffield{Statement} \emph{Types.} A restructuring changes the scope, assets or financing of a company:
\begin{itemize}
\item \emph{Investment actions} add businesses: acquisitions by merger, consolidation, tender offer or asset purchase; joint ventures; minority stakes \citep[pp.~702--704]{investment-valuation-2012}.
\item \emph{Divestment actions} remove them: sales to another company or to managers, spin-offs and split-offs that distribute a subsidiary's shares to shareholders, and equity carve-outs that sell part of a subsidiary to the public. Damodaran lists spin-offs and split-offs among the ways to return cash and value to shareholders \citep[pp.~437--438]{applied-corporate-finance-2014}.
\item \emph{Restructuring actions} change cost or capital structure: cost-reduction programmes and closures, leveraged recapitalisations, debt-for-equity swaps, sale-and-leasebacks, and reorganisation in insolvency.
\end{itemize}

\emph{Motives} \citep[pp.~706--710]{investment-valuation-2012}:
\begin{itemize}
\item \emph{Operating synergy}: economies of scale, pricing power, combined functional strengths, faster growth.
\item \emph{Financial synergy}: cash slack matched with projects, added debt capacity, tax benefits.
\item \emph{Undervaluation} of the target.
\item \emph{Control}: replacing poor management.
\item \emph{Diversification}: valuable only for owners who cannot diversify themselves.
\item \emph{Managerial self-interest}: empire building, pay and hubris (Roll 1986).
\end{itemize}
Divestments are motivated by focus, the conglomerate discount, a better owner for the unit, regulatory or tax reasons, or the need for cash. Diversified firms trade at lower values than portfolios of focused firms \citep[pp.~2--3]{lang-stulz-1993}. After hostile takeovers, about 60 percent of acquirers made significant divestitures of unrelated units, refocusing the target \citep[pp.~709--710]{investment-valuation-2012}.

\emph{Initial evaluation} (Interpretation). Before valuation, establish:
\begin{itemize}
\item what is proposed and why now;
\item whether it is material to the company's value, leverage or earnings;
\item the terms: price, premium, form of payment, financing, conditions and timetable;
\item the stated rationale against the evidence for that type of deal;
\item who bears the risk and who gains, including management's incentives.
\end{itemize}

\emph{Evidence.} Targets earn large returns: about 30 percent in tender offers and 20 percent in mergers (Jensen and Ruback 1983). Bidders earn little, and about half earn negative returns. Studies of post-deal performance find that many acquisitions fail to earn their cost of capital: KPMG found 17 percent created value and 53 percent destroyed it. A large share are later divested \citep[pp.~704--705, 723--724]{investment-valuation-2012}.

\emph{Evaluating each action.}
\begin{itemize}
\item \emph{Investment actions} add value only if the value of the target plus synergies exceeds the price, after integration costs. The synergies must be specific, valued with their timing and probability, and set against the premium.
\item \emph{Divestments} add value if the unit is worth more to a buyer, or as an independent company, than it contributes to the parent. Count tax leakage, dis-synergies such as lost shared services, and stranded central costs.
\item \emph{Cost restructurings} are capital projects: one-off costs now for savings later, valued by NPV with realistic savings and timing (\hyperref[meth:capital-project-evaluation]{projects}).
\item \emph{Balance-sheet restructurings} change value through taxes, distress costs, incentives and signals (\hyperref[mod:capital-structure-theory]{capital structure}). A leveraged recapitalisation raises the tax shield and discipline, and also distress risk and the cost of both capital sources.
\end{itemize}

\qffield{Practice} \emph{Original.}
\begin{itemize}
\item \emph{Cost restructuring.} A further plant programme at Harbour would cost 12 now, deductible, so 9 after tax, and save 8 a year before tax for five years. At 8 percent the NPV is $-9+6\times3.9927=14.96$. If the savings start a year late or reach only half, the NPV falls to 13.18 or 2.98. Harbour's 2025 restructuring charge of 18 is a reminder that costs can exceed the plan.
\item \emph{Spin-off.} A group trades at 1,500. Analysts value the parent without the unit at 1,200 and the unit independently at 450, so a spin-off would reveal 150, a 9.1 percent discount to the sum. If separation creates dis-synergies of 20 a year after tax, capitalised at 8 percent they are worth 250, which more than removes the gain. The spin-off should go ahead only if the dis-synergies can be avoided or are much smaller.
\item \emph{Acquisition.} Harbour's bolt-on is evaluated in \hyperref[ex:restructuring-valuation-effects]{restructuring effects}.
\end{itemize}

\qffield{Conditions} Evidence on acquisition returns is mostly from US deals in the 1970s to 1990s as summarised by Damodaran; later studies vary by period and deal type. Legal forms of mergers, spin-offs and tax-free reorganisations differ across jurisdictions, and the tax treatment can decide the form.

\qffield{Pitfalls} Accepting management's synergy estimates without timing, costs and probability. Valuing a divestment without stranded costs. Treating a restructuring charge as non-recurring when restructurings recur. Calling a deal good because it raises EPS.

\qffield{See also} \hyperref[ex:restructuring-valuation-effects]{valuation effects}; \hyperref[meth:valuation-approach-choice]{sum of the parts}; \hyperref[meth:corporate-governance-evaluation]{governance and acquisitions}.

\end{qfmethod}

\begin{qfexample}[Valuing a restructuring and its effects: Harbour's bolt-on acquisition valued by comparable companies, precedent transactions and DCF, value created against value transferred, and the effects on EPS, net debt to EBITDA and the WACC]\label{ex:restructuring-valuation-effects}
\qffield{Statement} \emph{Valuation methods} for a restructuring:
\begin{itemize}
\item \emph{Comparable companies}: trading multiples of listed peers. These give a minority, non-control value.
\item \emph{Comparable (precedent) transactions}: multiples paid in recent acquisitions, which include control premiums and synergies. They can be biased upwards because the deals that happened are those where bidders paid most \citep[p.~725]{investment-valuation-2012}.
\item \emph{Premiums paid}: for a listed target, the premium over its undisturbed price in comparable deals \citep[pp.~702--705]{investment-valuation-2012}.
\item \emph{Discounted cash flow}: the target valued on its own, then with synergies, each at the target's cost of capital, not the bidder's (p.~725).
\end{itemize}
The value of synergy is the value of the combined firm with synergy less the standalone values. Synergy should be valued with its form, its timing and its cost (pp.~715--718). In Damodaran's example, Compaq--Digital synergy of \$2,422 million would be worth \$1,724.86 million if delayed three years (pp.~717--718).

The \emph{value created} for the buyer is standalone value plus synergy less price; the rest of the synergy is \emph{transferred} to the seller through the premium. A control premium must not be counted twice, once in transaction multiples and again as a separate premium (p.~725).

\emph{Effects.}
\begin{itemize}
\item \emph{EPS.} A deal is accretive if the target's earnings exceed the after-tax cost of financing it. Damodaran calls the accretion test meaningless as a value test: a high-P/E acquirer buying a lower-P/E target with stock always shows accretion, yet its P/E falls (p.~724).
\item \emph{Leverage.} A debt-financed deal raises net debt to EBITDA, which matters for covenants and ratings.
\item \emph{Cost of capital.} The WACC changes with the debt ratio, the business mix and the rating, but the target's cash flows should be discounted at a rate for their own risk.
\end{itemize}

\qffield{Practice} \emph{Harbour Tools' bolt-on} (\hyperref[ex:integrated-credit-decision]{credit decision}). Harbour pays 60, financed with debt at 5.2 percent, for a distributor with revenue of 50 and EBITDA of 9 (6.67 times). Fair values: PP\&E 10 over five years, intangibles 20 over ten years, goodwill 30. Completion is 1 April 2026.

\emph{Valuation.}
\begin{itemize}
\item \emph{Comparable companies.} Listed distributors at 7.5 times EBITDA give 67.5.
\item \emph{Precedent transactions.} Recent deals at 8.0 times give 72.0, a figure that includes other buyers' synergies and premiums.
\item \emph{Premiums paid.} Not applicable to a private target.
\item \emph{DCF.} Harbour's DCF values the target standalone at 52, at the target's own cost of capital; integration synergies are worth 14 in present value after integration costs.
\end{itemize}
Harbour creates $52+14-60=6$, and transfers 8 of the 14 to the seller. The multiples suggest the price is fair to low. The DCF says the deal works only if the synergies are delivered: without them, Harbour pays 8 more than standalone value.

\emph{EPS} (full year, 2025 basis). The target's EBIT is $9-2-2=5.00$ after depreciation and amortisation of the fair-value uplifts. Interest at 5.2 percent on 60 is 3.12, so it adds $(5.00-3.12)\times0.75=1.41$ after tax. EPS rises from 0.5517 to 0.5658, by 2.5 percent.

\emph{Return on the price.} EBIT over price, $5.0/60=8.33$ percent, is a pre-tax return and cannot be compared with an after-tax WACC. After tax the deal earns $5.0\times0.75/60=6.25$ percent on the price in its first full year, 1.77 points below Harbour's 8.02 percent WACC and below the target's own cost of capital if its risk is similar. Three measures answer different questions:
\begin{itemize}
\item \emph{EPS accretion} (+2.5 percent) shows only that the after-tax return of 6.25 percent exceeds the after-tax cost of the debt, 3.90 percent; it ignores the cost of the equity risk the debt adds.
\item \emph{The accounting return} (6.25 percent, before synergies) is a one-year yield on the price; it says nothing about growth.
\item \emph{DCF value} compares the price with the present value of all future cash flows: standalone 52 against a price of 60, so $-8$ without synergies and $+6$ with them.
\end{itemize}
The first-year after-tax return below the WACC is consistent with the standalone DCF: without synergies the deal destroys value, whatever it does to EPS.

\emph{Leverage.} Net debt to EBITDA rises from 1.63 to 1.85 pro forma, with a full year of the target's EBITDA and the price added to net debt. That is within the 2.50 covenant but above the 1.75 dividend lock-up, so the dividend stops until leverage falls (\hyperref[ex:integrated-credit-decision]{option D}).

\emph{WACC.} Debt rises to 415 against equity, including the options, of 1,126.19, so D/E is 36.85 percent. The relevered beta is 1.1870 and the cost of equity 9.53 percent.
\begin{itemize}
\item If the cost of debt stayed at 5.2 percent, the WACC would be 8.01 percent, almost unchanged from 8.02.
\item On the illustrative rating schedule (\hyperref[meth:capital-structure-benchmarking]{benchmarking}), coverage of $116/(415\times k_d)$ falls to the BB+ band. The cost of debt becomes 6.0 percent and the WACC 8.17 percent.
\end{itemize}
The deal is not cheap capital. Its hurdle rate is the target's own cost of capital, not the change in Harbour's WACC.

\emph{Recommendation.} Proceed only with a synergy plan worth at least 8, the premium over standalone value, with milestones; keep the debt within the covenant and accept the dividend lock-up; and walk away above about 66. At 66 the 14 of synergy would be wholly transferred. A seller's request for a price near the precedent multiple of 72 should be refused. The recommendation rests on the DCF comparison, which the after-tax return confirms; the EPS accretion adds nothing to it.

\qffield{Conditions} The figures are hypothetical. The standalone value and synergy are inputs to this example. The rating schedule is illustrative, and the lender's facility terms in the case, not the schedule, set Harbour's actual borrowing cost.

\qffield{Pitfalls} Using transaction multiples as evidence of value to this buyer. Discounting target cash flows at the bidder's WACC. Treating EPS accretion as value creation. Comparing a pre-tax return on the price with an after-tax WACC. Counting the control premium twice. Ignoring that leverage after the deal changes the cost of both debt and equity.

\qffield{See also} \hyperref[meth:restructuring-evaluation]{evaluating restructurings}; \hyperref[meth:enterprise-value-multiples]{EV multiples}; \hyperref[meth:wacc-calculation]{WACC}; \hyperref[meth:intangible-acquisition-routes]{acquired intangibles}.

\end{qfexample}
